\documentclass[preprint,12pt,numbers]{elsarticle}

\usepackage{amsmath,amssymb,bm}
\usepackage{graphicx}
\usepackage{booktabs}
\usepackage{multirow}
\usepackage{threeparttable}
\usepackage{siunitx}
\usepackage{xcolor}
\usepackage{hyperref}
\usepackage{enumitem}
\usepackage{url}
\usepackage{float}
\usepackage{longtable}
\usepackage{array}

\journal{Automation in Construction}

\begin{document}
\begin{frontmatter}

\title{SAFE-3D: Site-Aligned Association and Fusion with Uncertainty Estimation for UAV-Referenced Multi-Camera Construction Safety Monitoring}

\author[skku]{Sujin Jin}
\author[skku]{Seunghee Park\corref{cor1}}
\cortext[cor1]{Corresponding author.}
\address[skku]{Department of Global Smart City, Sungkyunkwan University, Suwon, Republic of Korea}

\begin{abstract}
Two-dimensional detections from construction-site cameras are not directly comparable across viewpoints, and image-plane proximity is distorted by depth and perspective. This study proposes SAFE-3D, a UAV-referenced framework that converts worker and equipment detections from multiple fixed CCTV cameras into uncertainty-aware metric three-dimensional observations and global trajectories in a common site coordinate system. A UAV-reconstructed site map defines the world frame, each CCTV camera is registered through robust 2D--3D correspondence and perspective-$n$-point estimation, and detections are lifted into 3D using a representative ground-contact anchor and scale-corrected depth. Anisotropic observation covariances encode depth- and viewpoint-dependent uncertainty and are used consistently for multi-view fusion, Mahalanobis gating, Kalman filtering, and motion-gated tracklet re-linking. Across four complementary datasets, CCTV-to-map registration achieved a mean reprojection error of 4.6 pixels. Leave-one-out validation against ten total-station-measured worker locations yielded a 3D RMS error of 1.270 m. Localization ablation reduced reprojection error from 53.0 to 12.9 pixels through the successive use of the ground-contact anchor, metric scale correction, and covariance-aware fusion. The latest association configuration achieved an IDF1 of 0.845 and a HOTA score of 0.807 without scene-specific Euclidean distance threshold tuning, outperforming fixed-distance baselines with IDF1 scores of 0.424--0.473. On an operating construction site, worker and equipment IDF1 scores were 0.67 and 0.76, respectively. Image-plane proximity cues remained camera dependent, with AUC values of 0.770 and 0.629, and 24.6\% of visually close pairs in one camera lying outside the corresponding 3D proximity threshold. These results show that SAFE-3D extends existing CCTV infrastructure from view-specific tracking to metric, site-coordinate 3D safety monitoring.
\end{abstract}

\begin{keyword}
Construction safety \sep Multi-camera tracking \sep UAV mapping \sep 3D localization \sep Uncertainty-aware association \sep CCTV
\end{keyword}

\end{frontmatter}

\section{Introduction}
\label{sec:introduction}

Construction remains one of the most hazardous industries worldwide. The spatially dynamic nature of construction sites creates persistent risks because workers, heavy equipment, materials, temporary structures, and active work zones continuously change their positions and interactions. Effective safety monitoring must therefore extend beyond recognizing objects or unsafe behavior in isolated video frames and should support continuous interpretation of where workers and equipment are located, how they move, and whether their spatial relationships indicate hazardous proximity.

Positioning technologies such as GPS/GNSS, RFID, UWB, and BLE have been used to localize construction resources and maintain persistent identities. These technologies can provide direct positional measurements, but they require tags or devices to be attached to each tracked entity and introduce operational burdens associated with deployment, calibration, power, communication coverage, maintenance, and retrieval. Positional sensors also do not inherently provide visual information about personal protective equipment, surrounding work context, or observable worker behavior. For these reasons, considerable research has investigated non-contact construction monitoring using existing CCTV infrastructure and computer vision.

Recent advances in deep learning have enabled automated hardhat and harness detection, unsafe-behavior recognition, and worker and equipment tracking from construction videos. Pose estimation and person re-identification have also been introduced to reduce track fragmentation and identity switches when workers wear visually similar personal protective equipment. Nevertheless, image coordinates are not physical site coordinates. Two-dimensional observations from different cameras are governed by different projective geometries and cannot be compared directly. The same pixel separation can correspond to substantially different physical distances depending on object depth and camera viewpoint. Single-camera trajectories are further interrupted by occlusion, missed detections, illumination changes, and limited field of view.

Multi-camera methods address some of these limitations by connecting identities across views. Construction-oriented studies have used re-identification for safety-compliance monitoring, matched earthmoving equipment across non-overlapping cameras, and combined appearance, location, and motion cues for cross-camera worker association. BIM-integrated monitoring has also provided spatial context by visualizing camera-derived observations in three-dimensional building models. However, these approaches generally treat cross-camera identity matching and physical localization as separate problems. They do not consistently represent each camera observation as a metric 3D estimate with explicit uncertainty in a common construction-site coordinate system.

Three limitations therefore remain. First, most construction multi-camera tracking methods preserve identity continuity without transforming all observations into a common metric 3D site frame. Consequently, trajectories and inter-object distances are not consistently interpretable in physical units across cameras. Second, geometry-based approaches commonly rely on planar homographies, simplified ground-plane assumptions, or scene-specific Euclidean association thresholds. Such assumptions are restrictive in irregular and evolving construction environments, and they do not account explicitly for the fact that localization uncertainty varies with camera pose, object distance, view direction, occlusion, and depth estimation quality. Third, previous studies often evaluate detection, tracking, localization, or visualization separately. End-to-end evidence linking camera-to-map registration, metric 3D lifting, uncertainty-aware global association, field deployment, and safety-relevant spatial analysis remains limited.

To address these limitations, this study proposes SAFE-3D, standing for Site-Aligned Association and Fusion with Uncertainty Estimation in 3D. SAFE-3D uses a UAV-reconstructed site map as a common metric world frame for multiple fixed CCTV cameras. Each CCTV camera is registered to the site map, and each detection is lifted into metric 3D using a representative ground-contact anchor and scale-corrected depth. An anisotropic covariance is assigned to each observation to represent depth- and viewpoint-dependent uncertainty. The same uncertainty representation is then used for inverse-covariance fusion, Mahalanobis-distance gating, Kalman filtering, and motion-gated tracklet re-linking. The objective is not to introduce a new detector or single-camera tracker, but to transform existing camera-specific outputs into physically interpretable global trajectories in a shared construction-site coordinate system.

The contributions of this study correspond directly to the three limitations identified above:

\begin{enumerate}[leftmargin=*]
    \item \textbf{UAV-referenced metric 3D localization across multiple CCTV cameras.} A UAV-reconstructed site map establishes a common metric frame, and each fixed CCTV camera is registered to this map. Camera-specific detections are converted into site-coordinate 3D observations without requiring a separately hand-constructed ground-plane homography for every camera.

    \item \textbf{Uncertainty-aware multi-view fusion and global association.} Depth- and viewpoint-dependent anisotropic covariance is modeled during 2D-to-3D lifting. The covariance is used consistently for observation fusion, statistical gating, trajectory estimation, and tracklet re-linking, avoiding scene-specific Euclidean distance threshold tuning.

    \item \textbf{End-to-end validation from controlled experiments to operating construction sites.} Detection, camera registration, metric localization, component ablation, cross-camera association, trajectory reconstruction, and proximity analysis are evaluated across four complementary datasets. The analysis further quantifies the camera dependence and depth ambiguity of image-plane proximity cues.
\end{enumerate}

The remainder of this paper is organized as follows. Section~\ref{sec:related_work} reviews vision-based construction monitoring, sensor-based localization, multi-camera tracking, and spatial registration. Section~\ref{sec:method} presents SAFE-3D. Section~\ref{sec:results} reports the experimental results. Section~\ref{sec:discussion} discusses the findings, practical implications, and limitations. Section~\ref{sec:conclusion} concludes the paper.

\section{Related Work}
\label{sec:related_work}

\subsection{Vision-Based Construction Safety Monitoring}
Computer vision has been widely investigated as a non-contact complement to manual construction-safety inspection. Previous studies have reviewed object detection, posture estimation, activity analysis, and behavior-based safety monitoring in construction environments. Subsequent work has addressed hardhat non-compliance, safety-harness verification, personal protective equipment classification, and worker re-identification. These studies demonstrate that safety-relevant semantics can be extracted from site imagery. Their outputs, however, are generally image-level detections, classifications, behavioral labels, or camera-specific tracks rather than continuous trajectories in a common metric site coordinate system.

\subsection{Sensor-Based Localization of Construction Resources}
RFID, UWB, GPS/GNSS, and BLE have been applied to worker--equipment proximity monitoring and construction-resource localization. Sensor-based systems directly measure position and can maintain a unique identity for each instrumented resource. Their practical limitations include device installation, calibration, power management, communication coverage, maintenance, and retrieval. In addition, positional measurements alone do not describe PPE status or visual work context. These limitations motivate vision-based localization methods that can provide both semantic and spatial information without instrumenting every tracked entity.

\subsection{Multi-Camera Worker and Equipment Tracking}
Early construction studies tracked workers within individual views or detected workers and equipment from site video streams. Detector--tracker pipelines subsequently became common, but similar PPE, occlusion, missed detections, and abrupt viewpoint changes continue to cause trajectory fragmentation and identity switches. Pose-assisted re-identification has improved discrimination among workers wearing visually similar clothing. Construction-specific multi-camera studies have matched earthmoving equipment across non-overlapping cameras, accumulated PPE-compliance records across views, and combined appearance, location, and motion cues through message passing. These methods advance cross-camera identity continuity, but they generally do not represent every observation as a metric 3D position with an associated localization covariance in a shared site frame.

\subsection{Geometry-Based Association and Spatial Registration}
General multi-camera tracking research has used homographies, space--time--view graphs, lifted multicut formulations, and ground-coordinate Kalman filters to exploit geometric consistency. These studies show that geometry and uncertainty can resolve ambiguities that image-plane overlap or appearance features alone cannot. However, planar homographies assume that relevant objects can be represented on a common plane and often require manually selected correspondences for each camera. Such assumptions are restrictive in sites containing slopes, level changes, uneven terrain, temporary structures, and evolving work zones.

Within construction research, BIM and UAV mapping have been used to provide spatial context for camera observations. BIM-integrated approaches visualize detected workers in three-dimensional models, while UAV photogrammetry provides metric site reconstruction and supports spatial analysis. Ground-to-aerial localization methods register ground imagery to aerial references to estimate camera pose. Nevertheless, BIM-integrated monitoring commonly depends on an available model and limited camera configurations, whereas UAV mapping and cross-view localization mainly focus on map generation or camera localization. A unified framework that uses a UAV-derived metric site map as the common reference for multiple fixed CCTV cameras and propagates localization uncertainty through global association remains underdeveloped.

\subsection{Research Gap and Positioning}
SAFE-3D treats spatial localization and identity association as a single geometrically consistent tracking problem. The UAV-derived site map provides a common metric frame; camera-specific observations are lifted into 3D with uncertainty; and the same covariance representation supports fusion, statistical gating, trajectory estimation, and tracklet re-linking. This design differs from methods that perform appearance-based cross-camera matching first and attach spatial information afterward. It also differs from planar surveillance formulations by retaining full 3D site coordinates suitable for construction-safety distance analysis.

\begin{table*}[t]
\centering
\caption{Functional comparison of representative approaches and SAFE-3D.}
\label{tab:related_comparison}
\begin{tabular}{lccccc}
\toprule
Approach & Continuous tracking & Cross-camera association & Metric site frame & Explicit uncertainty & Safety output \\
\midrule
Pose-assisted ReID & Yes & No & No & No & Partial \\
ReID + PPE monitoring & Yes & Yes & No & No & Yes \\
Multi-camera equipment monitoring & Yes & Yes & No & No & No \\
BIM + computer vision & Yes & No & Yes & No & Yes \\
Message-passing worker tracking & Yes & Yes & No & No & Partial \\
Lifted geometric multi-camera tracking & Yes & Yes & Partial & No & No \\
Ground-coordinate uncertainty tracking & Yes & No & Partial & Yes & No \\
\textbf{SAFE-3D} & \textbf{Yes} & \textbf{Yes} & \textbf{Yes} & \textbf{Yes} & \textbf{Yes} \\
\bottomrule
\end{tabular}
\end{table*}

\section{Methodology}
\label{sec:method}

\subsection{Framework Overview}
SAFE-3D comprises an offline geometry stage and an online tracking stage. The offline stage reconstructs a UAV site map, aligns it to a metric world frame, and registers each fixed CCTV camera to the map. The online stage detects workers and equipment, extracts a representative ground-contact anchor, lifts each observation into metric 3D, estimates its anisotropic covariance, and associates observations with predicted global tracks. Multi-view observations assigned to the same track are fused before Kalman updating. A motion-gated tracklet re-linking stage recovers identity continuity after longer observation gaps.

\subsection{UAV Site Reconstruction and Metric Alignment}
UAV imagery is reconstructed using sparse structure-from-motion camera poses and a dense reconstruction model. Let the reconstruction coordinate system be $\mathcal{M}$ and the metric site frame be $\mathcal{W}$. A similarity transformation is estimated from metric control correspondences:
\begin{equation}
(s_g,\mathbf{R}_g,\mathbf{t}_g)=\arg\min_{s,\mathbf{R},\mathbf{t}}\sum_i\left\|\bar{\mathbf{X}}_i^{W}-(s\mathbf{R}\bar{\mathbf{X}}_i^{M}+\mathbf{t})\right\|_2^2.
\end{equation}
The transformation is applied consistently to UAV camera poses, sparse and dense points, and depth-derived geometry. The resulting map provides a common metric frame, reference UAV views for CCTV registration, and a pixel-to-3D lookup from reference imagery to map points.

\subsection{CCTV-to-Map Registration}
The intrinsic matrix $\mathbf{K}_c$ of each CCTV camera $c$ is obtained through calibration and lens distortion is removed. Dense correspondences are extracted between a CCTV reference frame and overlapping UAV reference images. A UAV reference pixel is mapped to its 3D site point, producing 2D--3D correspondences $\{(\mathbf{x}_{i}^{c},\mathbf{X}_{i}^{W})\}$. Camera extrinsics are initialized with PnP-RANSAC and refined by robust reprojection minimization:
\begin{equation}
(\mathbf{R}_c,\mathbf{t}_c)=\arg\min_{\mathbf{R},\mathbf{t}}\sum_{i\in\mathcal{I}_c}\rho\!\left(\left\|\pi\!\left(\mathbf{K}_c(\mathbf{R}\mathbf{X}_i^W+\mathbf{t})\right)-\mathbf{x}_i^c\right\|_2\right).
\end{equation}
Registration quality is checked using the inlier count, inlier ratio, median reprojection error, and positive-depth consistency.

\subsection{Detection and Ground-Contact Anchor}
Workers and equipment are detected independently in each camera. Association is permitted only within the same object class. For a bounding box $b=(x_1,y_1,x_2,y_2)$, SAFE-3D defines a representative ground-contact anchor (GCA)
\begin{equation}
u_g=\frac{x_1+x_2}{2},\qquad v_g=y_1+\alpha_g(y_2-y_1),
\end{equation}
where $\alpha_g$ is fixed across datasets. The GCA is not claimed to be a direct foot or tire-contact detector; it is a stabilized image anchor intended to represent the object's ground location more consistently than the box center or raw bottom-center.

\subsection{Depth-Scale Correction and 2D-to-3D Lifting}
Relative depth $D_c(\mathbf{x})$ is converted to metric camera depth through a camera-specific scale estimated from CCTV--map correspondences. For valid correspondence samples,
\begin{equation}
s_c=\arg\min_s\sum_j\rho\left(z_{j}^{\mathrm{map}}-sD_c(\mathbf{x}_j)\right)^2.
\end{equation}
The metric depth at the GCA is $z_g=s_cD_c(\mathbf{a})$. The camera-frame point is
\begin{equation}
\mathbf{X}^{c}=z_g\mathbf{K}_c^{-1}[u_g,v_g,1]^\top,
\end{equation}
and the world-frame observation is obtained using the inverse camera extrinsic transform.

\subsection{Observation Uncertainty}
Depth error is generally larger along the viewing ray than perpendicular to it and increases with object distance. SAFE-3D therefore assigns each observation an anisotropic covariance in the camera frame, with a larger variance along the ray direction. The covariance is rotated into the world frame and augmented with terms representing anchor instability and camera registration uncertainty. This covariance is a geometry-informed uncertainty approximation used for relative weighting and compatibility testing; it is not presented as a fully calibrated probabilistic guarantee.

\subsection{Uncertainty-Aware Global Association and Fusion}
Each global track follows a constant-velocity Kalman model. For predicted track $k$ and observation $i$, the innovation and combined covariance are
\begin{equation}
\mathbf{r}_{ki}=\mathbf{z}_i-\mathbf{H}\hat{\mathbf{x}}_k^{-},\qquad
\mathbf{S}_{ki}=\mathbf{H}\mathbf{P}_k^{-}\mathbf{H}^{\top}+\boldsymbol{\Sigma}_i.
\end{equation}
The squared Mahalanobis distance is
\begin{equation}
d_{ki}^{2}=\mathbf{r}_{ki}^{\top}\mathbf{S}_{ki}^{-1}\mathbf{r}_{ki}.
\end{equation}
Pairs satisfying $d_{ki}^{2}\leq\chi^2_{3,0.95}$ are retained and a one-to-one assignment is obtained with the Hungarian algorithm. The confidence level, process noise, confirmation length, and retention period are fixed system parameters shared across datasets. Thus, the method avoids scene-specific Euclidean distance threshold tuning but is not described as parameter-free.

When multiple cameras assign observations to the same track at the same time, they are fused using inverse-covariance weighting:
\begin{equation}
\bar{\boldsymbol{\Sigma}}=\left(\sum_i\boldsymbol{\Sigma}_i^{-1}\right)^{-1},\qquad
\bar{\mathbf{z}}=\bar{\boldsymbol{\Sigma}}\sum_i\boldsymbol{\Sigma}_i^{-1}\mathbf{z}_i.
\end{equation}
The fused observation and covariance are then used for the Kalman update.

\subsection{Motion-Gated Tracklet Re-Linking}
The conservative online gate avoids erroneous merges but can fragment a true identity after prolonged occlusion or missed detection. A post-tracking re-linking stage considers terminated and subsequently initialized tracklets only when they share the same class, do not overlap temporally, and are separated by less than a fixed maximum gap. The terminal state of the earlier tracklet is propagated to the start time of the later tracklet using the same constant-velocity model. Compatibility is evaluated by the covariance-normalized distance
\begin{equation}
d_{\mathrm{link}}^2=(\mathbf{p}^{+}-\hat{\mathbf{p}}^{-})^{\top}(\mathbf{P}^{-}+\mathbf{P}^{+})^{-1}(\mathbf{p}^{+}-\hat{\mathbf{p}}^{-}).
\end{equation}
Pairs satisfying the same $\chi^2$ criterion are matched one-to-one with the Hungarian algorithm. Re-linking recovers fragmentation without relaxing the online gate or introducing a site-specific Euclidean radius.

\subsection{Safety-Relevant Trajectory Outputs}
Each confirmed global track stores time, position, velocity, covariance, and object class. Pairwise distance is calculated between representative ground-contact anchors in the metric world frame. Distances to polygonal or three-dimensional hazard regions are computed directly from the reconstructed trajectories. These outputs provide a geometric basis for safety analysis; they do not by themselves constitute a complete warning policy, which would additionally require equipment envelopes, speed, operational rules, and risk-dependent thresholds.

\section{Results}
\label{sec:results}

\subsection{Experimental Design and Datasets}
SAFE-3D was evaluated on four complementary datasets to separate component-level validation from end-to-end field evaluation. Dataset A comprised four CCTV cameras in a controlled indoor environment and was used for detection, localization consistency, component ablation, and cross-camera association. Dataset B combined UAV imagery, three fixed CCTV cameras, and total-station measurements in an outdoor controlled environment and was used for CCTV-to-map registration, metric localization validation, proximity analysis, and outdoor association. Dataset C was collected at an operating construction site using a UAV-derived site map and three CCTV cameras and was used to evaluate worker and equipment association and site-coordinate trajectory reconstruction. Dataset D contained three CCTV cameras at a separate construction site without a UAV map and was used only as an auxiliary geometry-controlled evaluation of the association module and its sensitivity to missed detections. Dataset D was not treated as an end-to-end validation of the UAV-referenced pipeline.

The manually annotated evaluation subset of Dataset C contained 10 worker identities and seven equipment identities. Dataset A contained two worker identities and was intended to isolate method differences rather than establish statistical generalization. Dataset B involved three workers. Dataset D consisted of two short independent sessions. Camera layout, frame counts, synchronization, and annotation procedures should be reported in the Appendix once the final acquisition log is confirmed.

\begin{table}[t]
\centering
\caption{Role of the four evaluation datasets.}
\label{tab:datasets}
\begin{tabular}{p{0.8cm}p{3.4cm}p{3.6cm}}
\toprule
Set & Environment and configuration & Primary purpose \\
\midrule
A & Controlled indoor, four CCTV cameras & Detection, localization, ablation, association \\
B & Controlled outdoor, UAV + three CCTV cameras & Registration, metric validation, proximity, association \\
C & Operating site, UAV + three CCTV cameras & Field association and trajectory reconstruction \\
D & Separate site, three CCTV cameras & Auxiliary association and recall sensitivity \\
\bottomrule
\end{tabular}
\end{table}

\subsection{Object Detection Performance}
Detection was evaluated as the upstream observation source for localization and association rather than as the main methodological contribution. Datasets A and B used a COCO-pretrained YOLO11 detector with the person class retained. Datasets C and D used a detector retrained for worker and construction-equipment classes because of the domain gap associated with small distant objects, partial occlusion, and construction-specific machinery.

Dataset A achieved an overall precision of 0.945, recall of 0.937, F1 score of 0.941, mAP@0.5 of 0.944, and mAP@0.5:0.95 of 0.936. Dataset B achieved an overall precision of 0.977, recall of 0.956, F1 score of 0.966, mAP@0.5 of 0.990, and mAP@0.5:0.95 of 0.740. The lower high-IoU mAP in Dataset B reflects outdoor scale variation despite stable person detection.

On Dataset C, the field-trained detector achieved an overall precision of 0.899, recall of 0.793, F1 score of 0.842, mAP@0.5 of 0.870, and mAP@0.5:0.95 of 0.592. Worker and equipment mAP@0.5 values were 0.852 and 0.889, respectively. Dataset D achieved an overall precision of 0.838, recall of 0.787, F1 score of 0.811, mAP@0.5 of 0.818, and mAP@0.5:0.95 of 0.497. Worker and equipment mAP@0.5 values were 0.871 and 0.765. The controlled-to-field performance gap reflects the visual complexity of real sites and explains why long-distance false negatives and partial occlusion remain major sources of trajectory interruption.

\begin{table*}[t]
\centering
\caption{Object detection performance across datasets.}
\label{tab:detection}
\begin{tabular}{llccccc}
\toprule
Dataset & Class & Precision & Recall & F1 & mAP@0.5 & mAP@0.5:0.95 \\
\midrule
A & Worker & 0.945 & 0.937 & 0.941 & 0.944 & 0.936 \\
B & Worker & 0.977 & 0.956 & 0.966 & 0.990 & 0.740 \\
C & Worker & 0.902 & 0.737 & 0.811 & 0.852 & 0.583 \\
C & Equipment & 0.895 & 0.849 & 0.871 & 0.889 & 0.600 \\
D & Worker & 0.884 & 0.842 & 0.862 & 0.871 & 0.563 \\
D & Equipment & 0.792 & 0.731 & 0.760 & 0.765 & 0.431 \\
\bottomrule
\end{tabular}
\end{table*}

\subsection{Metric 3D Localization}
\subsubsection{Multi-View Consistency and Total-Station Validation}
Dataset A evaluated internal cross-view consistency. For synchronized observations of the same worker, the distance from each camera-specific 3D estimate to the cross-view centroid had a median of 12.2 cm, an RMSE of 17.7 cm, and a 90th percentile of 29.4 cm. The median cross-view reprojection error was 12.9 pixels. Metric scale was independently checked using four measured ArUco-marker separations spanning approximately 2--6 m, and reconstructed distances remained within \(\pm3\%\). The exact marker IDs, measurement instrument, and pairwise values should be listed in the Appendix.

Dataset B provided external validation. The three CCTV cameras were registered to the UAV-derived site map with reprojection errors of 4.2--4.8 pixels and a mean of 4.6 pixels. Ten total-station-measured worker ground locations were used to evaluate the 3D localization pipeline. A similarity transform estimated and evaluated on all ten points produced an in-sample RMS residual of 0.776 m, which characterizes alignment fit rather than independent localization accuracy. In leave-one-out validation, the transform was estimated from nine points and evaluated on the excluded point, yielding an RMS error of 1.270 m. This leave-one-out value is used as the primary external localization metric. Camera-to-centroid consistency in Dataset B had a median of 21.5 cm and an RMSE of 33.8 cm, reflecting the increased viewing distances and outdoor geometry.

\begin{table}[t]
\centering
\caption{Metric localization validation.}
\label{tab:localization}
\begin{tabular}{p{4.4cm}p{2.4cm}}
\toprule
Metric & Result \\
\midrule
Dataset A camera-to-centroid median / RMSE / P90 & 12.2 / 17.7 / 29.4 cm \\
Dataset A cross-view reprojection median & 12.9 px \\
Dataset A metric-scale error & within \(\pm3\%\) \\
Dataset B CCTV-to-map reprojection & 4.2--4.8 px (mean 4.6) \\
Dataset B ten-point in-sample residual & 0.776 m \\
Dataset B leave-one-out RMS error & 1.270 m \\
Dataset B camera-to-centroid median / RMSE & 21.5 / 33.8 cm \\
\bottomrule
\end{tabular}
\end{table}

Dataset C was collected during ongoing site operations. Installing and repeatedly surveying total-station targets throughout active work zones was not feasible without interfering with construction activities. Dataset C was therefore evaluated through cross-camera identity consistency and trajectory continuity rather than independent point-wise absolute localization. Dataset D used a ground-plane mapping and was excluded from metric 3D localization evaluation.

\subsubsection{Localization Ablation}
Ablation was conducted on 1,132 matched observations in Dataset A. The box-center configuration produced a median reprojection error of 53.0 pixels and pairwise consistency error of 40.5 cm. Moving the anchor to the raw box bottom-center reduced these values to 49.5 pixels and 35.5 cm. The proposed GCA further reduced them to 30.6 pixels and 24.7 cm. Adding metric depth-scale correction reduced the errors to 16.9 pixels and 16.3 cm. Finally, covariance-aware fusion reduced fused-estimate reprojection error to 12.9 pixels. Pairwise consistency remained 16.3 cm because it was computed from camera-specific observations before fusion.

\begin{table}[t]
\centering
\caption{Localization ablation on Dataset A.}
\label{tab:ablation}
\begin{tabular}{lcccc}
\toprule
Variant & Anchor & Scale & Fusion & Reproj. / Consist. \\
\midrule
V1 & Box center & No & No & 53.0 px / 40.5 cm \\
V2 & Bottom-center & No & No & 49.5 px / 35.5 cm \\
V3 & GCA & No & No & 30.6 px / 24.7 cm \\
V4 & GCA & Yes & No & 16.9 px / 16.3 cm \\
SAFE-3D & GCA & Yes & Yes & 12.9 px / 16.3 cm \\
\bottomrule
\end{tabular}
\end{table}

\subsection{Cross-Camera Global Association}
The association evaluation examined whether covariance-normalized matching could outperform fixed-distance alternatives while using the same detections, timestamps, classes, and spatial observations. No-association evaluation on Dataset A produced an IDF1 of 0.263 because each camera retained independent local identities. Fixed-radius 3D nearest matching and Euclidean--Hungarian assignment achieved IDF1 values of 0.424 and 0.473, respectively. The latest SAFE-3D configuration achieved an IDF1 of 0.845 and a HOTA score of 0.807 without scene-specific Euclidean distance threshold tuning. A GT-tuned Euclidean reference achieved an IDF1 of 0.982, but this value is treated only as an oracle upper reference because its optimal radius was selected using the evaluation identities.

The practical advantage of SAFE-3D is not the absence of all parameters. Rather, the same statistical confidence level and system parameters are shared across datasets, whereas a fixed Euclidean radius changes with camera layout, object distance distribution, and density. The radius sensitivity analysis should therefore be presented as evidence of site-specific threshold dependence rather than as a claim of parameter-free operation.

The motion-gated re-linking stage was introduced to recover identities after longer observation gaps. Older gate-only values and draft re-linking values were removed from the principal comparison to avoid mixing configurations. The principal manuscript reports the latest complete SAFE-3D run. Any additional IDSW, Pair-F1, or reconstructed-ID values should be inserted only when they are confirmed to originate from that same run.

Dataset C yielded IDF1 scores of 0.67 for workers and 0.76 for equipment. The higher equipment score is consistent with larger image extent, more stable detections, and slower motion. The evaluation subset produced 13 reconstructed worker identities for 10 GT identities and six reconstructed equipment identities for seven GT identities. Worker over-fragmentation was concentrated around long occlusions; the missing equipment identity should be attributed to a confirmed log-based failure mode before final submission rather than to an unverified explanation.

Dataset D is retained as an auxiliary experiment that isolates association from UAV reconstruction and CCTV-to-map registration errors. Its ground mapping was established from GT correspondences, so its result should be interpreted as association behavior under a controlled alternative geometry rather than as evidence that the full SAFE-3D pipeline is independent of UAV mapping. The measured overall cross-camera Pair-F1 across the two sessions was 0.452, compared with 0.315 for a frame-wise greedy baseline and 0.504 for the GT-box geometric upper reference.

\begin{table}[t]
\centering
\caption{Confirmed association results used in the principal manuscript.}
\label{tab:association}
\begin{tabular}{p{3.9cm}cc}
\toprule
Method / target & IDF1 & HOTA \\
\midrule
Dataset A: no association & 0.263 & -- \\
Dataset A: fixed-radius 3D nearest & 0.424 & -- \\
Dataset A: Euclidean--Hungarian & 0.473 & -- \\
Dataset A: SAFE-3D, latest run & 0.845 & 0.807 \\
Dataset A: GT-tuned Euclidean oracle & 0.982 & -- \\
Dataset C: workers & 0.67 & -- \\
Dataset C: equipment & 0.76 & -- \\
\bottomrule
\end{tabular}
\end{table}

\subsection{End-to-End Pipeline and Safety-Relevant Outputs}
Dataset B demonstrated the complete UAV--CCTV geometry pipeline. Eighty sampled UAV images were used for dense reconstruction with sparse SfM poses, three CCTV cameras were registered to the map, and reconstructed detections were expressed in the same metric site frame. The use of 80 sampled images should be described as a computationally constrained subset selected from the full UAV acquisition; the final Appendix should report the total acquired image count and sampling rule. The pipeline achieved 80.9\% trajectory coverage in the evaluated interval.

Dataset C demonstrated field trajectory reconstruction for workers and equipment on the UAV/BEV site map. This result extends camera-specific detections to global site-coordinate trajectories, but the manuscript should avoid claiming absolute field localization accuracy because no independent total-station validation was available in this operating-site dataset.

For Dataset B, 2,651 synchronized worker pairs were analyzed. The median estimated separation between representative ground-contact anchors was 2.12 m. The minimum draft value of 0.11 m was removed from the principal text because the anchor-to-anchor estimate should not be interpreted as physical clearance between object boundaries and because threshold-near events require uncertainty-aware interpretation.

Image-plane distance was evaluated against proximity labels derived from metric 3D trajectories. Pixel-distance AUC was 0.770 for Camera 1 and 0.629 for Camera 3, with best F1 scores of 0.724 and 0.677. In Camera 1, 24.6\% of pairs judged close in image space were outside the corresponding 3D proximity threshold. The previous 3D AUC row was removed; the 3D distance provides the reference spatial relation in this comparison rather than a competing classifier.

\begin{table}[t]
\centering
\caption{Camera dependence of 2D proximity cues relative to metric 3D proximity labels.}
\label{tab:proximity}
\begin{tabular}{lccc}
\toprule
Camera & AUC & Best F1 & False-close rate \\
\midrule
Camera 1 & 0.770 & 0.724 & 24.6\% \\
Camera 3 & 0.629 & 0.677 & -- \\
\bottomrule
\end{tabular}
\end{table}

\section{Discussion}
\label{sec:discussion}

\subsection{Principal Findings}
SAFE-3D extends camera-specific construction monitoring into a common site-coordinate representation by combining a UAV-derived 3D map with multiple fixed CCTV cameras. The results support three main findings. First, the UAV map provides a geometric reference that allows observations from otherwise incomparable camera views to be interpreted in one metric frame. Second, the localization ablation shows that a representative ground-contact anchor, metric depth-scale correction, and covariance-aware fusion contribute sequentially to geometric consistency. Third, the reconstructed trajectories support physically interpretable spatial outputs, including worker--worker separation, worker--equipment relations, and distances to predefined hazard regions.

The validation results should be interpreted at the appropriate level. Dataset A demonstrates strong internal cross-view consistency and isolates the contribution of each localization component. Dataset B provides the principal independent metric validation through leave-one-out evaluation against total-station measurements. Dataset C demonstrates field-level identity continuity and trajectory reconstruction but does not establish independent absolute point accuracy. Dataset D isolates the association module under an alternative ground mapping and is therefore auxiliary rather than an end-to-end validation of the UAV-referenced pipeline.

\subsection{Implications for Construction Safety Monitoring}
Vision-based construction-safety research has predominantly advanced image-space recognition of PPE, unsafe behavior, and interactions. However, collision, intrusion, and equipment-approach risks are governed by physical relationships rather than visual separation alone. SAFE-3D reinterprets existing CCTV cameras as spatially aware sensors by registering them to a metric site frame. This connection is particularly relevant because UAV mapping, as-built documentation, progress monitoring, and digital-twin workflows already generate static spatial representations, whereas CCTV provides persistent observations of dynamic site actors. SAFE-3D links these static and dynamic information sources in one coordinate system.

The proximity analysis quantifies why this connection matters. Pixel-distance discrimination varied substantially between cameras and produced false-close cases caused by depth ambiguity. A single image-plane threshold therefore cannot provide a camera-independent safety distance. Metric 3D trajectories do not automatically constitute a complete warning system, but they provide the spatial substrate required to define equipment envelopes, restricted regions, and distance-based rules in physical units.

\subsection{Role of Ground-Contact Localization and Uncertainty}
The box center does not represent the ground position of a worker or item of equipment. The raw bottom-center is closer to the expected contact region but remains sensitive to bounding-box jitter, shadows, truncation, and partial lower-body occlusion. The proposed GCA is therefore described conservatively as a representative ground anchor rather than a directly detected footpoint. Its role is to reduce systematic image-anchor inconsistency before depth-based lifting.

Scale correction is equally important because relative monocular depth cannot be interpreted directly in meters. Aligning depth-derived geometry to the metric site frame produced the largest reduction after anchor correction in the ablation study. The final fusion stage further reduced reprojection error by weighting camera observations according to their uncertainty.

The uncertainty model is most valuable when observations differ in quality. A nearby frontal observation, a distant oblique observation, and a partially occluded observation should not contribute equally to fusion or association. Observation covariance allows SAFE-3D to weight these cases differently and to interpret spatial discrepancy relative to predicted uncertainty. However, the present covariance should be viewed as a geometry-informed approximation. A stronger future evaluation should examine calibration using empirical coverage, normalized innovation statistics, or error stratification by distance and viewing direction.

\subsection{Cross-Camera Association and Threshold Transfer}
The association results show that a fixed Euclidean radius is sensitive to site geometry. A radius that is sufficiently large to retain distant uncertain observations can merge nearby objects, whereas a smaller radius can reject correct matches under long viewing distances. The GT-tuned Euclidean result demonstrates that distance-based association can perform well when the evaluation identities are available for tuning, but this condition is not representative of a newly installed site.

SAFE-3D does not eliminate all parameters. It replaces scene-specific Euclidean radius tuning with a statistical compatibility rule based on propagated track covariance and observation covariance. The confidence level, process noise, confirmation length, and retention settings remain fixed system parameters. This distinction is important: the contribution is cross-site reuse of a common uncertainty-aware association rule, not parameter-free tracking.

The re-linking stage addresses a complementary failure mode. Conservative online gating favors track fragmentation over incorrect merging during long observation gaps. Motion-gated re-linking recovers compatible tracklets afterward while retaining the same covariance-normalized criterion. It should therefore be interpreted as a continuity-recovery mechanism rather than an independent tracking contribution. A complete final table should report IDF1, HOTA, Pair-F1, IDSW, and reconstructed identity count from the same latest run to quantify this trade-off without mixing experimental configurations.

\subsection{Practical Deployment Considerations}
The system naturally separates into offline and online stages. UAV reconstruction, metric alignment, and CCTV registration can be performed offline, whereas detection, lifting, association, fusion, and trajectory update operate on incoming camera streams. Construction-site geometry is not static, however. Excavation, backfilling, structural growth, and temporary facilities can invalidate parts of the map. Practical deployment therefore requires a map-update policy based on periodic UAV acquisition or change detection.

CCTV pose accuracy directly limits localization accuracy. Cameras displaced by vibration, impact, or maintenance require re-registration. Long-term deployment would benefit from automatic pose-drift detection using repeated reference matching. Operational visualization should expose global tracks, camera support, uncertainty, and zone distance on a site or BEV map so that managers can interpret why a spatial relation was generated.

Downstream warning logic should combine distance with equipment type, operating radius, speed, work-zone rules, and uncertainty. Near a decision threshold, a probabilistic or interval-based policy is more defensible than a deterministic threshold applied to a single estimated distance. The current study establishes the geometry and trajectory representation needed for such rules but does not claim to validate a complete hazard-warning policy.

\subsection{Limitations and Future Work}
Several limitations remain. First, localization accuracy depends on CCTV calibration, map quality, depth estimation, and GCA reliability. Distant workers, prolonged occlusion, and incomplete lower-body observations increase uncertainty and may require semantic ground constraints, stronger temporal smoothing, or explicit contact-point estimation. Second, the operating-site dataset lacks independent point-wise total-station validation because continuous target installation and surveying were not feasible during active operations. Future field studies should use non-interfering survey protocols, robotic total stations, or synchronized reference sensors.

Third, cross-camera association remains sensitive to missed detections, severe occlusion, synchronization error, and dense interactions. Geometry and motion can be complemented by short-term appearance evidence, although visually similar PPE limits its role. Fourth, the four datasets cover complementary conditions but remain limited in site type, camera height, equipment diversity, density, illumination, and construction phase. Broader multi-site validation is required. Fifth, the map-maintenance interval and the effect of progressive site change were not evaluated. Future work should investigate dynamic map updates, online pose refinement, calibrated uncertainty, uncertainty-aware risk rules, and integration with BIM and construction digital twins.

\section{Conclusion}
\label{sec:conclusion}

This study proposed SAFE-3D, a UAV-referenced framework for site-aligned association and fusion with uncertainty estimation in multi-camera construction monitoring. The framework registers fixed CCTV cameras to a UAV-derived metric site map, lifts representative ground-contact observations into 3D using scale-corrected depth, models anisotropic localization uncertainty, and constructs global worker and equipment trajectories through covariance-aware fusion, Mahalanobis gating, Kalman filtering, and motion-gated tracklet re-linking.

Across controlled, metric-reference, and operating-site datasets, the framework achieved a mean CCTV-to-map reprojection error of 4.6 pixels and a leave-one-out 3D localization RMS error of 1.270 m against ten total-station measurements. Localization ablation reduced reprojection error from 53.0 to 12.9 pixels through the successive use of the ground-contact anchor, metric scale correction, and covariance-aware fusion. The latest cross-camera association configuration achieved an IDF1 of 0.845 and a HOTA score of 0.807 without scene-specific Euclidean distance threshold tuning, compared with IDF1 scores of 0.424--0.473 for fixed-distance baselines. At an operating construction site, worker and equipment IDF1 scores were 0.67 and 0.76, respectively. The proximity analysis further showed that image-plane distance is camera dependent and susceptible to depth-induced false-close relations.

These findings demonstrate that existing CCTV infrastructure can be extended from view-specific detection and tracking to metric, site-coordinate 3D trajectory monitoring. SAFE-3D provides a geometric and uncertainty-aware basis for subsequent worker--equipment proximity analysis, restricted-zone monitoring, and digital-twin integration. Future work will focus on dynamic site-map maintenance, online camera-pose refinement, calibrated uncertainty, broader multi-site validation, and risk rules that combine distance, speed, equipment envelope, and localization confidence.


\section*{Acknowledgements}
Funding information will be inserted after confirmation of the applicable National Research Foundation of Korea grant number.

\section*{CRediT authorship contribution statement}
\textbf{Sujin Jin}: Conceptualization, Methodology, Software, Validation, Formal analysis, Investigation, Data curation, Visualization, Writing--original draft. \textbf{Seunghee Park}: Conceptualization, Supervision, Project administration, Funding acquisition, Writing--review and editing.

\section*{Declaration of competing interest}
The authors declare that they have no known competing financial interests or personal relationships that could have appeared to influence the work reported in this paper.

\section*{Data availability}
The data that support the findings of this study are subject to construction-site access and privacy restrictions. Derived evaluation annotations and implementation details may be made available from the corresponding author upon reasonable request, subject to approval by the relevant site stakeholders.

\bibliographystyle{elsarticle-num}
\bibliography{refs}

\end{document}
